\subsection{Écrans coordonnés et dualité des rayons réciproques}
\label{exc:pantallas-dualidad}

La réduction sélectionnée par \(K\) et la construction de Leech
proviennent de lectures différentes de la même incidence. Pour préciser
leur coordination, on réunit ici les applications algébriques qui conservent
ces deux données. L'espace de rang onze et le réseau de rang vingt-quatre
ne sont pas identifiés ; on spécifie leurs quotients respectifs et
l'involution qui agit sur le facteur de routes et de mémoire.

\subsubsection{Résidu, quotient et changement de section}

L'évaluation arithmétique relevée d'APP conserve
\(n=9Q+r\), avec \(r\in\{1,\ldots,9\}\) dans la section
résiduelle positive. Le transport TPK conserve séparément ces
coordonnées, leur orientation et la mémoire de leurs changements. Pour
représenter un échange entre les deux coordonnées, il faut
un domaine stable par échange. On utilise donc l'espace ambiant
\(\ZZ^2\), dont les entiers signés prolongent les coordonnées de
la coordinatisation ; cela n'affirme pas que toutes ses paires soient des trajectoires
admissibles d'un même historique TPK.

Sur \(\mathscr D_0=\ZZ^2\times\RR_{>0}\), la réalisation quadratique
déclarée est
\begin{equation}
 q_R(m,w)=\frac{m^2}{R^2}+w^2R^2,
 \qquad \mathsf T_0(m,w,R)=(w,m,R^{-1}).
 \label{eq:dualidad-normalizada-k}
\end{equation}
La variable \(R\) est une échelle sans dimension de la coordinatisation de lecture.
Les symboles \(m,w\) désignent ici les deux coordonnées entières
transportées. Leur reconnaissance ultérieure comme impulsion et enroulement
maintient cette distinction entre construction et réalisation.

\begin{proposition}[Échange et forme quadratique]
\label{prop:dualidad-radios-k}
L'application \(\mathsf T_0\) est une involution et satisfait
\[
 q_{R^{-1}}(w,m)=q_R(m,w).
\]
Sur \(\ell^2(\ZZ^2)\), l'opérateur unitaire
\(U_T|m,w\rangle=|w,m\rangle\) entrelace les opérateurs
diagonaux \(H(R)|m,w\rangle=q_R(m,w)|m,w\rangle\) :
\[
 U_TH(R)U_T^{-1}=H(R^{-1}).
\]
\end{proposition}
\begin{proof}
Deux échanges restituent les coordonnées et deux inversions restituent
\(R\). La substitution donne
\(q_{R^{-1}}(w,m)=w^2R^2+m^2/R^2\). L'échange permute
la base orthonormée ; il est donc unitaire, et cette égalité vérifie
l'entrelacement sur chaque vecteur de la base. Celui-ci s'étend aux
domaines des opérateurs diagonaux par la même égalité des poids.
\end{proof}

La section résiduelle importe. Pour \(n=9\), les couples
\((r,Q)=(9,0)\) et \((0,1)\) représentent le même entier ;
leurs évaluations \(81/R^2\) et \(R^2\) ne coïncident pas pour tout
\(R\). Le changement de section est incorporé exactement au moyen de
\[
 L_9=\ZZ(9,-1),\qquad S(x_1,x_2)=(x_2,x_1).
\]
Définissons, sur les classes du quotient,
\begin{equation}
 \mathcal E_R^{L_9}([x])=\min_{\ell\in L_9}q_R(x+\ell),
 \qquad
 \mathsf T_{\mathrm{cov}}([x]_{L_9},R)
       =([Sx]_{SL_9},R^{-1}).
 \label{eq:dualidad-cociente-k}
\end{equation}

\begin{proposition}[Dualité compatible avec le changement de section]
Les deux expressions de \eqref{eq:dualidad-cociente-k} sont bien
définies. Les sections positive et nulle, avec l'augmentation
correspondante du quotient, représentent la même classe. De plus,
\[
 \mathcal E_{R^{-1}}^{SL_9}([Sx])=\mathcal E_R^{L_9}([x]).
\]
\end{proposition}
\begin{proof}
Remplacer \(x\) par \(x+\ell_0\) permute les candidats au minimum.
Celui-ci existe parce que \(q_R(x+k(9,-1))\) est une fonction quadratique coercive
en \(k\in\ZZ\). Le changement \((9,Q)\mapsto(0,Q+1)\)
a pour différence \((9,-1)\). Enfin,
\[
 \min_{\ell\in L_9}q_{R^{-1}}(Sx+S\ell)
 =\min_{\ell\in L_9}q_R(x+\ell).
\]
Puisque \(S^2=I\), l'échange inverse restitue la classe et l'échelle.
\end{proof}

L'égalité dépend du transport également de la relation d'équivalence,
de \(L_9\) à \(SL_9\). On conserve ainsi l'information du changement
de section qu'un simple remplacement du représentant \(9\) par
\(0\) ferait perdre.

\subsubsection{L'écran réticulaire et sa coordination avec \(K\)}

Soit \(\Lambda=N_\alpha\), le réseau de Leech construit dans
\ref{exc:leech}. Ajoutons le plan hyperbolique entier
\(U=\ZZ e_+\oplus\ZZ e_-\), avec
\(e_+^2=e_-^2=0\) et \(\langle e_+,e_-\rangle=1\).
Alors \(L_{26}=\Lambda\oplus U\) a pour signature \((25,1)\) et
\[
 e_+^\perp=\Lambda\oplus\ZZ e_+,
 \qquad e_+^\perp/\ZZ e_+\cong\Lambda.
\]
En effet, si \(y=\lambda+ae_++be_-\), la condition
\(\langle y,e_+\rangle=0\) équivaut à \(b=0\), et le quotient
élimine exactement la composante \(ae_+\). Son application est
\(q_{24}(\lambda+ae_+)=\lambda\).

Pour la réduction \(P_{10}\) déjà démontrée, soit
\[
 \mathscr G_K=\{(v_{11},P_{10}v_{11}):v_{11}\in V_{11}\}.
\]
On définit
\begin{equation}
 \mathcal R_K(v,d,y)
 =\bigl((P_{11}v,P_{10}v),d,q_{24}(y)\bigr),
 \label{eq:k-mapa-pantallas}
\end{equation}
de domaine \(\RR^{12}\times\mathscr D_0\times e_+^\perp\)
et de codomaine \(\mathscr G_K\times\mathscr D_0\times\Lambda\).

\begin{theorem}[Isomorphisme des écrans quotientés]
\label{thm:k-pantallas-coordinadas}
L'application \eqref{eq:k-mapa-pantallas} induit un isomorphisme
\[
 (\RR^{12}/\RR\mathbf1)\times\mathscr D_0
       \times(e_+^\perp/\ZZ e_+)
 \ \cong\ \mathscr G_K\times\mathscr D_0\times\Lambda.
\]
Elle entrelace les transformations qui agissent par \(\mathsf T_0\)
sur le facteur \(\mathscr D_0\) et par l'identité sur les autres.
\end{theorem}
\begin{proof}
\(P_{11}\) induit un isomorphisme de \(\RR^{12}/\RR\mathbf1\)
sur \(V_{11}\). Puisque \(P_{10}P_{11}=P_{10}\), l'application
\([v]\mapsto(P_{11}v,P_{10}v)\) identifie ce quotient au
graphe \(\mathscr G_K\) ; l'inverse prend la première composante
et sa classe. Le quotient réticulaire vient d'être identifié par
\(q_{24}\). Le facteur restant est conservé par l'identité.
Le produit des trois applications établit l'isomorphisme.
L'involution et les projecteurs agissent sur des facteurs distincts ;
leurs compositions commutent donc. La forme quadratique est conservée
par la proposition \ref{prop:dualidad-radios-k}.
\end{proof}

Le graphe retient \(v_{11}\) avec sa projection. Cet
isomorphisme n'affirme donc pas que \(V_{11}\) et \(V_{10}\) sont isomorphes :
la réduction isolée a pour noyau \(\ell_K\), tandis que le graphe
conserve la coordonnée antérieure. Cette précision identifie ce que conserve
chaque application et permet de relier les écrans \(11\to10\) et
\(26\to24\) sans les confondre.

Dans leur ensemble, \(K\) détermine la direction de réduction ; le drapeau
dodécaphasique détermine le voisin réticulaire utilisé pour Moonshine ;
et l'échange des coordonnées de route et de mémoire conserve la
forme quadratique sous des rayons réciproques. L'application à une théorie
physique des cordes ou à un espace de dimension onze requiert en outre
son action, ses champs et ses opérateurs de réalisation. Le résultat réuni
ici est l'application algébrique avec ses preuves, qui constitue un
antécédent explicite de ce développement ultérieur.
